\thispagestyle{fancy}
\fancyhead[R]{A.U 2025-2026}
\fancyhead[L]{ chapitre 6: Release 4 : Marketing et fidélisation }
\chapter{Release 4 : Marketing et fidélisation}
\section*{Introduction}
Vendre ne suffit pas : un commerçant doit aussi attirer de nouveaux clients et faire revenir les anciens. Ce chapitre présente la quatrième release, qui lui en donne les moyens. Le Sprint 7 couvre les promotions, les bannières ciblées et l'e-mailing ; le Sprint 8 couvre le programme de fidélité à points et ses niveaux.

%==================================================================
\section{Backlog du Sprint 7}
\Needspace{18\baselineskip}
Le tableau~\ref{tab:backlog-sprint7} présente le \textbf{backlog du Sprint 7}.

\renewcommand{\arraystretch}{1.05}
\setlength{\tabcolsep}{3pt}
\begin{longtable}{|>{\raggedright\arraybackslash}p{4cm}|>{\raggedright\arraybackslash}p{9cm}|>{\centering\arraybackslash}p{1.2cm}|}
\hline
\textbf{User Story} & \textbf{Tâches} & \textbf{Estimation} \\
\hline
\endfirsthead
\hline
\textbf{User Story} & \textbf{Tâches} & \textbf{Estimation} \\
\hline
\endhead
\hline
\multicolumn{3}{r}{\textit{Suite à la page suivante}} \\
\endfoot
\hline
\endlastfoot

\textbf{En tant que commerçant, je souhaite créer des codes promo}
& Développer l'entité \texttt{Coupon} : type (pourcentage, montant fixe, livraison offerte, cadeau, « un acheté, un offert »), valeur, montant minimum, période et plage horaire, limites globale et par client, segment, catégories et produits ciblés. & 1 \\
\cline{2-3}
& Développer l'API CRUD, l'activation / désactivation et le suivi des utilisations. & 1 \\
\cline{2-3}
& Calculer les statistiques par code (commandes, revenus, taux de conversion). & 1 \\
\cline{2-3}
& Développer l'interface « Promotions ». & 1 \\
\cline{2-3}
& Tester la fonctionnalité. & 1 \\
\hline

\textbf{En tant que commerçant, je souhaite créer des remises automatiques}
& Développer l'API des remises sur un produit ou une catégorie, avec période de validité. & 1 \\
\cline{2-3}
& Afficher le prix barré et le badge de remise sur la vitrine. & 1 \\
\cline{2-3}
& Tester la fonctionnalité. & 1 \\
\hline

\textbf{En tant que commerçant, je souhaite publier des bannières ciblées}
& Développer l'API des bannières : titre, images ordinateur et mobile, vidéo, bouton d'action, alignement, animation, durée, planning. & 1 \\
\cline{2-3}
& Cibler une bannière par segment de clientèle et par type d'appareil. & 1 \\
\cline{2-3}
& Développer le formulaire de bannière avec aperçu ordinateur et mobile. & 1 \\
\cline{2-3}
& Afficher les bannières dans le carrousel d'accueil de chaque gabarit. & 1 \\
\hline

\textbf{En tant que commerçant, je souhaite envoyer des e-mails à mes clients}
& Intégrer l'API d'envoi transactionnel de Brevo (clé en variable d'environnement). & 1 \\
\cline{2-3}
& Développer l'envoi d'une newsletter à un segment et le contact individuel d'un client. & 1 \\
\cline{2-3}
& Développer l'interface « E-mail marketing » et ses statistiques. & 1 \\
\hline

\multicolumn{2}{|r|}{\textbf{Total}} & \textbf{15} \\
\hline
\caption{Backlog du Sprint 7}
\label{tab:backlog-sprint7}
\end{longtable}

\section{Diagramme de cas d'utilisation du Sprint 7}
\Needspace{20\baselineskip}
\begin{figure}[H]
    \centering
    \img{0.8}{cu_sprint7.png}
    \caption{Diagramme de cas d'utilisation du Sprint 7}
    \label{fig:cu-sprint7}
\end{figure}

\section{Services web du Sprint 7}
Le tableau~\ref{tab:ws-sprint7} présente les services web du Sprint 7.
\renewcommand{\arraystretch}{1.15}
\setlength{\tabcolsep}{4pt}
\begin{longtable}{|p{1.5cm}|p{6.8cm}|p{6.2cm}|}
\hline
\textbf{Méthode} & \textbf{URL} & \textbf{Description} \\
\hline
\endfirsthead
\hline
\textbf{Méthode} & \textbf{URL} & \textbf{Description} \\
\hline
\endhead
\multicolumn{3}{|c|}{\textbf{Promotions --- /api/v1/admin/promotions (order-service)}} \\
\hline
GET & /api/v1/admin/promotions/stats & Statistiques des promotions. \\
\hline
GET, POST & /api/v1/admin/promotions/coupons & Lister, créer un code promo. \\
\hline
PUT, DELETE & /api/v1/admin/promotions/coupons/\{id\} & Modifier, supprimer un code promo. \\
\hline
PATCH & /api/v1/admin/promotions/coupons/\{id\}/toggle & Activer ou désactiver un code. \\
\hline
GET, POST & /api/v1/admin/promotions/discounts & Lister, créer une remise. \\
\hline
PUT, DELETE & /api/v1/admin/promotions/discounts/\{id\} & Modifier, supprimer une remise. \\
\hline
PATCH & /api/v1/admin/promotions/discounts/\{id\}/toggle & Activer ou désactiver une remise. \\
\hline
\multicolumn{3}{|c|}{\textbf{Bannières --- marketing-service}} \\
\hline
GET, POST & /api/v1/admin/banners & Lister, créer une bannière. \\
\hline
PUT, DELETE & /api/v1/admin/banners/\{id\} & Modifier, supprimer une bannière. \\
\hline
PATCH & /api/v1/admin/banners/\{id\}/toggle & Activer ou désactiver une bannière. \\
\hline
GET & /api/v1/public/banners?position=\&segment= & Bannières actives pour une position et un segment. \\
\hline
\multicolumn{3}{|c|}{\textbf{E-mail marketing --- /api/v1/admin/email}} \\
\hline
GET & /api/v1/admin/email/stats & Statistiques d'envoi. \\
\hline
POST & /api/v1/admin/email/newsletter & Envoyer une newsletter. \\
\hline
POST & /api/v1/admin/email/contact/\{userId\} & Contacter un client. \\
\hline
\caption{Services web du Sprint 7}
\label{tab:ws-sprint7}
\end{longtable}

\section{Les cas d'utilisation du Sprint 7}
\subsection{Cas d'utilisation « Créer un code promo ciblé »}
\subsubsection{Description textuelle}
Le tableau~\ref{tab:uc-coupon} présente la description textuelle de ce cas.
\renewcommand{\arraystretch}{1.5}
\begin{longtable}{|p{4cm}|p{11cm}|}
\hline
\textbf{Acteurs} & Commerçant \\
\hline
\textbf{Objectif} & Créer une offre réservée à une cible et en mesurer l'efficacité. \\
\hline
\textbf{Pré-condition} & Le commerçant est connecté et sa boutique est active. \\
\hline
\textbf{Scénario principal} &
1. Le commerçant ouvre « Promotions » puis « Nouveau code ».\newline
2. Il saisit le code (par exemple \texttt{VIP20}), le type et la valeur de la remise, et le montant minimum de commande.\newline
3. Il définit la période, éventuellement une plage horaire, et les limites d'utilisation.\newline
4. Il restreint le code à un segment (par exemple VIP) et, s'il le souhaite, à des catégories ou des produits.\newline
5. Le système vérifie que le code n'existe pas déjà, calcule son statut (programmé, actif ou expiré) à partir des dates, et l'enregistre.\newline
6. Lorsqu'un client l'utilise, le système enregistre l'utilisation et met à jour le nombre de commandes, les revenus et le taux de conversion du code. \\
\hline
\textbf{Scénario alternatif} &
\textbf{Code déjà existant :} à l'étape 5, le système demande un autre code.\newline
\textbf{Limite atteinte :} à l'étape 6, si le client a déjà utilisé le code le nombre maximal de fois autorisé, le système refuse la remise avec un message explicite. \\
\hline
\caption{Description textuelle du cas « Créer un code promo ciblé »}
\label{tab:uc-coupon}
\end{longtable}

\subsubsection{Diagramme de séquence système}
\Needspace{18\baselineskip}
\begin{figure}[H]
    \centering
    \img{0.65}{dss_coupon.png}
    \caption{Diagramme de séquence système du cas « Créer un code promo ciblé »}
    \label{fig:dss-coupon}
\end{figure}

\subsubsection{Diagramme de séquence objet}
\Needspace{20\baselineskip}
\begin{figure}[H]
    \centering
    \img{0.9}{dso_coupon.png}
    \caption{Diagramme de séquence objet du cas « Créer un code promo ciblé »}
    \label{fig:dso-coupon}
\end{figure}

\section{Réalisation du Sprint 7}
\Needspace{20\baselineskip}
\textbf{Promotions :} la figure~\ref{fig:ui-promotions} présente la gestion des codes promo et des remises, avec leurs indicateurs.
\begin{figure}[H]
    \centering
    \img{0.85}{ui_promotions.png}
    \caption{Gestion des promotions}
    \label{fig:ui-promotions}
\end{figure}

\Needspace{20\baselineskip}
\textbf{Bannières :} la figure~\ref{fig:ui-bannieres} présente le formulaire de bannière, avec aperçu sur ordinateur et sur mobile.
\begin{figure}[H]
    \centering
    \img{0.85}{ui_bannieres.png}
    \caption{Création d'une bannière ciblée}
    \label{fig:ui-bannieres}
\end{figure}

\Needspace{20\baselineskip}
\textbf{E-mail marketing :} la figure~\ref{fig:ui-email} présente l'envoi d'une newsletter.
\begin{figure}[H]
    \centering
    \img{0.85}{ui_email_marketing.png}
    \caption{E-mail marketing}
    \label{fig:ui-email}
\end{figure}

%==================================================================
\section{Backlog du Sprint 8}
\Needspace{18\baselineskip}
Le tableau~\ref{tab:backlog-sprint8} présente le \textbf{backlog du Sprint 8}, consacré au programme de fidélité.

\renewcommand{\arraystretch}{1.05}
\setlength{\tabcolsep}{3pt}
\begin{longtable}{|>{\raggedright\arraybackslash}p{4cm}|>{\raggedright\arraybackslash}p{9cm}|>{\centering\arraybackslash}p{1.2cm}|}
\hline
\textbf{User Story} & \textbf{Tâches} & \textbf{Estimation} \\
\hline
\endfirsthead
\hline
\textbf{User Story} & \textbf{Tâches} & \textbf{Estimation} \\
\hline
\endhead
\hline
\multicolumn{3}{r}{\textit{Suite à la page suivante}} \\
\endfoot
\hline
\endlastfoot

\textbf{En tant que commerçant, je souhaite configurer le gain de points}
& Développer l'entité \texttt{LoyaltyConfig} (points par dinar, bienvenue, avis, anniversaire, promotion automatique). & 1 \\
\cline{2-3}
& Attribuer les points à la livraison d'une commande, à l'inscription et à la publication d'un avis, avec le multiplicateur du niveau. & 1 \\
\cline{2-3}
& Journaliser chaque mouvement de points (\texttt{PointsTransaction}). & 1 \\
\cline{2-3}
& Développer l'onglet « Configuration ». & 1 \\
\hline

\textbf{En tant que commerçant, je souhaite gérer les niveaux et ajuster les points d'un client}
& Développer l'API CRUD des segments (seuil de points, multiplicateur, avantages). & 1 \\
\cline{2-3}
& Développer la montée de niveau automatique après chaque attribution de points. & 1 \\
\cline{2-3}
& Développer l'ajustement manuel (le solde ne peut pas devenir négatif) et le classement des clients. & 1 \\
\cline{2-3}
& Développer les onglets « Avantages / Niveaux », « Classement » et « Ajustement manuel ». & 1 \\
\cline{2-3}
& Tester la fonctionnalité. & 1 \\
\hline

\textbf{En tant que client, je souhaite consulter mes points et mon niveau}
& Développer l'API \texttt{/profile/loyalty} (solde, niveau, historique, progression vers le niveau suivant). & 1 \\
\cline{2-3}
& Afficher la carte de fidélité dans « Mon profil ». & 1 \\
\hline

\multicolumn{2}{|r|}{\textbf{Total}} & \textbf{11} \\
\hline
\caption{Backlog du Sprint 8}
\label{tab:backlog-sprint8}
\end{longtable}

\section{Diagramme de cas d'utilisation du Sprint 8}
\Needspace{20\baselineskip}
\begin{figure}[H]
    \centering
    \img{0.8}{cu_sprint8.png}
    \caption{Diagramme de cas d'utilisation du Sprint 8}
    \label{fig:cu-sprint8}
\end{figure}

\section{Services web du Sprint 8}
Le tableau~\ref{tab:ws-sprint8} présente les services web du Sprint 8.
\renewcommand{\arraystretch}{1.15}
\setlength{\tabcolsep}{4pt}
\begin{longtable}{|p{1.5cm}|p{6.8cm}|p{6.2cm}|}
\hline
\textbf{Méthode} & \textbf{URL} & \textbf{Description} \\
\hline
\endfirsthead
\hline
\textbf{Méthode} & \textbf{URL} & \textbf{Description} \\
\hline
\endhead
\multicolumn{3}{|c|}{\textbf{Fidélité --- /api/v1/admin/loyalty (auth-service)}} \\
\hline
GET & /api/v1/admin/loyalty/config & Consulter la configuration du programme. \\
\hline
PUT & /api/v1/admin/loyalty/config & Modifier la configuration. \\
\hline
GET & /api/v1/admin/loyalty/leaderboard & Classement des clients par points. \\
\hline
POST & /api/v1/admin/loyalty/users/\{userId\}/adjust-points & Ajuster les points d'un client (positif ou négatif). \\
\hline
\multicolumn{3}{|c|}{\textbf{Segments --- /api/v1/admin/segments}} \\
\hline
GET, POST & /api/v1/admin/segments & Lister, créer un niveau. \\
\hline
PUT, DELETE & /api/v1/admin/segments/\{id\} & Modifier, supprimer un niveau. \\
\hline
GET & /api/v1/admin/users/by-segment/\{segmentName\} & Clients d'un niveau. \\
\hline
\multicolumn{3}{|c|}{\textbf{Espace client}} \\
\hline
GET & /api/v1/profile/loyalty & Solde, niveau et historique du client connecté. \\
\hline
\caption{Services web du Sprint 8}
\label{tab:ws-sprint8}
\end{longtable}

\section{Les cas d'utilisation du Sprint 8}
\subsection{Cas d'utilisation « Gagner des points et changer de niveau »}
\subsubsection{Description textuelle}
Le tableau~\ref{tab:uc-fidelite} présente la description textuelle de ce cas.
\renewcommand{\arraystretch}{1.5}
\begin{longtable}{|p{4cm}|p{11cm}|}
\hline
\textbf{Acteurs} & Client, Commerçant \\
\hline
\textbf{Objectif} & Récompenser automatiquement les achats et faire progresser le client dans les niveaux du programme. \\
\hline
\textbf{Pré-condition} & Le client est inscrit ; le programme de fidélité est configuré. \\
\hline
\textbf{Scénario principal} &
1. Le commerçant passe une commande du client au statut « livrée ».\newline
2. Le système calcule les points : $\text{points} = \text{arrondi}(\text{montant} \times \text{points par dinar} \times \text{multiplicateur du niveau})$. Par exemple, 120 DT avec 1 point par dinar et un multiplicateur de 1,5 donnent 180 points.\newline
3. Le système crédite le solde et journalise le mouvement.\newline
4. Si la promotion automatique est active, le système cherche le niveau le plus élevé dont le seuil est atteint et y fait passer le client.\newline
5. Le client consulte son solde, son niveau et sa progression dans « Mon profil ». \\
\hline
\textbf{Scénario alternatif} &
\textbf{Ajustement manuel :} le commerçant retire des points ; le solde est borné à zéro.\newline
\textbf{Pas de rétrogradation automatique :} une baisse du solde ne fait pas redescendre le client ; le niveau « Inactif » n'est attribué que manuellement. \\
\hline
\caption{Description textuelle du cas « Gagner des points et changer de niveau »}
\label{tab:uc-fidelite}
\end{longtable}

\subsubsection{Diagramme de séquence système}
\Needspace{18\baselineskip}
\begin{figure}[H]
    \centering
    \img{0.65}{dss_fidelite.png}
    \caption{Diagramme de séquence système du cas « Gagner des points et changer de niveau »}
    \label{fig:dss-fidelite}
\end{figure}

\subsubsection{Diagramme de séquence objet}
\Needspace{20\baselineskip}
\begin{figure}[H]
    \centering
    \img{0.9}{dso_fidelite.png}
    \caption{Diagramme de séquence objet du cas « Gagner des points et changer de niveau »}
    \label{fig:dso-fidelite}
\end{figure}

\section{Réalisation du Sprint 8}
\Needspace{20\baselineskip}
\textbf{Programme de fidélité :} les figures~\ref{fig:ui-fidelite-config} et~\ref{fig:ui-fidelite-niveaux} présentent la configuration du programme et la gestion des niveaux (Nouveau, Fidèle, VIP, Inactif).
\begin{figure}[H]
    \centering
    \img{0.85}{ui_fidelite_config.png}
    \caption{Configuration du programme de fidélité}
    \label{fig:ui-fidelite-config}
\end{figure}
\begin{figure}[H]
    \centering
    \img{0.85}{ui_fidelite_niveaux.png}
    \caption{Niveaux et avantages du programme de fidélité}
    \label{fig:ui-fidelite-niveaux}
\end{figure}

\Needspace{20\baselineskip}
\textbf{Espace client :} la figure~\ref{fig:ui-fidelite-client} présente la carte de fidélité du client dans son profil.
\begin{figure}[H]
    \centering
    \img{0.75}{ui_fidelite_client.png}
    \caption{Carte de fidélité dans l'espace client}
    \label{fig:ui-fidelite-client}
\end{figure}

\subsection*{Conclusion}
Cette release a donné au commerçant les leviers classiques du marketing en ligne : promotions ciblées, bannières, e-mails et fidélité. Ces outils reposent toutefois sur des segments définis à la main (Nouveau, Fidèle, VIP). Les deux dernières releases vont plus loin grâce à l'intelligence artificielle : l'essayage virtuel d'abord, puis l'analyse automatique du comportement des visiteurs.
